\section{Race of the Soul}

\begin{quotex}
There are many cases of people who are exactly of the same race of the body, of the same tribe, sometimes even, brothers or fathers and sons, of the same blood in the most real sense, but who nevertheless fail to understand each other. A barrier separates their souls, their way of feeling and seeing is different, and the common race of the body and common blood can do nothing about that. A possibility of comprehension exists, and therefore of true solidarity, of deep unity, only where a common ``race of the soul'' exists.
\end{quotex}


Since man is an ensouled body, his outward physical characteristics, or race of the body, are insufficient to determine him. That is, the body is neither a chunk of matter, nor an animal, so DNA will fail to define him. Keep in mind that by ``race'', \textbf{Julius Evola} does not mean how it is understood in the USA, with its emotionally charged history of slavery, apartheid, and civil rights. Rather, it is much more specific, so he can speak of a Nordic race, an Italic race, and so on. This is more obvious in languages like Italian or German, since the word for ``race'' in those languages is also used to refer to ``breeds'' of animals.

Applying phenomenology to the experience of race, we see, as did Evola, that sharp and deep divisions separate even those of the same physical race. These are not merely accidental qualities, but are constitutive of the person. These divisions often make mutual intelligibility impossible, without the effort to transcend such differences on a higher plane. There are two all too obvious consequences to this. First of all, any attempt to appeal to some racial solidarity of the first level will necessarily fail. Secondly, these divisions at the soul level will not be resolved by discussions or ``debates''. \textbf{Fabre d'Olivet} confirms this:

\begin{quotex}
Nature does not make men equal; souls differ still more than bodies.
\end{quotex}


From the esoteric perspective, what is of interest is why people hold certain views, not the views themselves. For example, a convinced materialist or atheist holds his view with a passionate intensity, since he is unable to conceive of, or experience interiorly, any other possibility that makes sense to him. Bracketing out for the moment the possibility of a spiritual or intellectual conversion, people remain locked into certain fixed patterns of soul life.

Evola also provides a different understanding of ``purity''. Physical purity is impossible, or at least, undesirable. There is a small, but not negligible, mutation rate in any population. Since most mutations are deleterious, an endogamous population will decline over the centuries. Purity, then, can only mean the adequation of corporeal and psychic qualities. Thus for example, on the one hand he can praise certain Nordic qualities, while also mentioning the degeneration of the Nordic race. Note that these qualities are all on the plane of the soul and spirit, not the body.

The takeaway is that successful movements depend on a core group of ``like-minded'' people, if that is your goal; great leaders instinctively understand this. As Evola explains,

\begin{quotex}
A possibility of comprehension, and therefore of true solidarity, of deep unity, exists only where a common ``race of the soul'' exists.
\end{quotex}


For example, desirable qualities may include ``calm and dominating superiority, solarity, sense of distance, active detachment''. Those that have such qualities actively can then identify those who may have them virtually. However, an indiscriminate mixing, a ``big tent'', won't come to much beyond putting some people in positions of power for its own sake. The telltale sign of degeneration is the failure to differentiate, which is also the hallmark of the Kali Yuga. Differentiation at the level of the soul is more important than the body. Evola concludes:

\begin{quotex}
it may happen that a person gives us the clear impression by his way of acting that he ``is of another race'' and, then, there is no longer anything to do with him; relationships of a various nature can exist with him, but always with an inner reservation and inner distance. He ``is no longer ours''.
\end{quotex}


A leader needs to make that determination. Of course, the most fundamental differentiation is that between a man and a woman, who relate to each other not just physically, but more importantly at the soul and spirit level. Nowadays, even that differentiation is losing force. Many, today, even fail to differentiate between humans and animals. The text follows.

Let's now return to the clarification of the three levels of the doctrine of race. We consider a theory of the race of the soul and a typology of the soul of the races as the doctrine of the second level. Such a doctrine needs to identify the elements, in their primary and irreducible modes, that act from the inside, causing groups of individuals to manifest a consistent mode of being or ``style'' as far as acting, thinking, and feeling are concerned. With this we arrive at a new concept of racial purity of a given type: it is no longer about, as in racial doctrine of the first level, seeing if a given individual manifests a specific group of physical characteristics or even character traits, that makes him conform to the inherited type, but it is about establishing if a given individual's race of the body is the adequate expression or confirmation of his race of the soul, and vice versa. If that is confirmed, the type is pure also in the sense of the investigations of the second level.

That then complements the results of the of the first level, because it no longer considers the various corporeal characteristics in the abstract, in a simple classification, which could even be disguised, instead of in living faces and individuals. It seeks instead to understand their secret, that is to say, what they express, the function for which they are used and by way of which, from case to case, they might even mean something different. As we already noted, a nose of a given form and an elongated, dolicephalous skull can meet both in a type of the races derived from the Nordic stock, or in an exemplar of the African race: but it is obvious that the two cases do not have the same significance. Besides, it is quite possible that a given type has predominant characteristics, e.g., Mediterranean in the anthropological race of the body -- so much so that the racial doctrine of the first level would assign him precisely to the Mediterranean race of the man of the West: nevertheless the latest research may notice that these Mediterranean traits, in the type in question, are assumed in a different function from what would normally be expected. The type that we are talking about uses them instead to express a soul, an internal attitude that is not Mediterranean, but, for example, Nordic or Near Eastern. That gives an absolutely different expressive value to the same traits and leads sometimes to certain distortions or alterations of the exterior Mediterranean element that are almost imperceptible in the research of the first level or are considered irrelevant and unimportant, while for the research of the second level they represent many ways to understand the ``inner race''. Here physiognomy, the study of the meaning of the human face, has an important role: it will however develop in a different direction from the preceding, which conceived each individual separately instead of as a member of a given superbiological community, of a given race of the soul.

On this higher plane, anthropology and paleoanthropology become valuable aids for the research into the original racial elements that entered into composition, were superimposed, or clashed at the dawn of civilizations. For the highest tasks of the race theory, it is insufficient to have considered the presence, for example in Italic origins, of a given number of typical skeletons and skulls and, integrating such researches with the archeological ones, to be able to affirm fundamentally the existence of an ancient pure Nordico-Aryan Italic human type. That would not get beyond an area of a museum. It is necessary, furthermore, to make this type speak, to penetrate into what a given corporeal form expresses, or what a given human structure is the symbol of. This is impossible without taking into account the field of the racial doctrine of the second level, and even to a certain extent the third level, disciplines that work with other methods of research and utilize another order of documents and evidence.

We can consider the so-called \emph{Rassenseelekunde} or psychoanthropology of L. F. Clauss as racial theories of the second level, insofar as it concerns his methods and general criteria. The necessity of such research was made clear by Clauss with convincing examples. For example, he considered the phenomena of comprehension. Actually, there are too many cases of people who are exactly of the same race of the body, of the same tribe, sometimes even, brothers or fathers and sons, of the same blood in the most real sense, but who nevertheless fail to understand each other. A barrier separates their souls, their way of feeling and seeing is different, and the common race of the body and common blood can do nothing about that.

A possibility of comprehension, and therefore of true solidarity, of deep unity, exists only where a common ``race of the soul'' exists. Subtle elements of an instinctive sensitivity come into question at this point. While for many years nothing was suspected, in a given circumstance it may happen that a person gives us the clear impression by his way of acting that he ``is of another race'' and, then, there is no longer anything to do with him; relationships of a various nature can exist with him, but always with an inner reservation and inner distance. He ``is no longer ours''. Usually, they spoke of character in this case. The expression is vague. There is not, in fact, a ``character'' in general but there are different modes of appearing of the dispositions of character conditioned by the inner race. For example, the way of being ``faithful'' of a Near Eastern people is different from that of a man of the Nordic or Dinaric race. The way of conceiving heroism by a Mediterranean man is different from that of a Japanese or a Russian, to use generic expressions and while not getting into precise classifications inherent to a doctrine of the race of the soul.

\flrightit{Posted on 2015-05-10 by Aeneas}

\begin{center}* * *\end{center}

\begin{footnotesize}
\begin{sffamily}
\texttt{Alistair Fraser on 2015-05-11 at 06:38 said:}

``it may happen that a person gives us the clear impression by his way of acting that he ``is of another race'' and, then, there is no longer anything to do with him; relationships of a various nature can exist with him, but always with an inner reservation and inner distance. He ``is no longer ours''. ''

Is it possible that by the power of early conditioning , the impulse that exists to express the genuine race of ones soul can be hindered greatly , so some people appear to be one type but are actually a repressed other type imprisoned withing their societal cage , but yet even as their true expression has been overwritten, blocked and given little room to sprout, still there may be signs of it that seep out through gesture and energy presence perceivable by others of similar ilk and so the ``virtually quality '' confirms

\hfill

\texttt{Paulo on 2015-05-11 at 15:49 said:}

well being born and raised in brasil,the eye of the hurricane kali yuga

i cant fail to notice the effects of this kind of mix

mix of the body,mix of the soul and mix of the sexes really,destiny had put me in a position were the chalenges will not give up on me!!!

\end{sffamily}
\end{footnotesize}

